% File: Tables/table_nominal_battery_summary.tex
% Usage: \NominalBatterySummaryTable[0.78\textwidth]
% Values obtained from the corrected nominal-case battery export.

\newcommand{\NominalBatterySummaryTable}[1][0.78\textwidth]{%
\begin{table}[H]
    \centering
    \caption{Initial state and main battery results for the nominal transient
    case using the selected \(6S1P\) configuration.}
    \label{tab:nominal_battery_summary}
    \footnotesize
    \renewcommand{\arraystretch}{1.18}
    \setlength{\tabcolsep}{6pt}
    \begin{tabularx}{#1}{
        |>{\raggedright\arraybackslash}X
        |>{\centering\arraybackslash}c
        |>{\centering\arraybackslash}c|
    }
        \hline
        \textbf{Item} &
        \textbf{Value} &
        \textbf{Unit} \\ \hline
        Initial state of charge
        & 60.00 & \% \\ \hline
        Final state of charge
        & 94.43 & \% \\ \hline
        Stored energy at the final instant
        & 62.889 & Wh \\ \hline
        Terminal-voltage range
        & 22.648--24.376 & V \\ \hline
        Maximum charge current
        & 1.500 & A \\ \hline
        Maximum discharge current
        & 0.946 & A \\ \hline
        Maximum charge power
        & 35.350 & W \\ \hline
        Maximum discharge power
        & 21.429 & W \\ \hline
        Battery-temperature range
        & 18.400--20.900 & \(^{\circ}\mathrm{C}\) \\ \hline
        Minimum charge-derating factor
        & 1.000 & -- \\ \hline
        Curtailed solar energy
        & 276.228 & Wh \\ \hline
        Unserved-load energy
        & \(\approx 0\) & Wh \\ \hline
        Integrated internal resistive loss
        & 0.252 & Wh \\ \hline
        Integrated conversion loss
        & 0.923 & Wh \\ \hline
    \end{tabularx}
\end{table}%
}
